\chapter*{\center \Large Abstract}

\noindent
Trouver le plus court trajet entre deux points d'une ville est une opération que
chacun utilise au quotidien, sans en soupçonner la difficulté : sur une carte
réelle, qui compte des dizaines à des centaines de milliers d'intersections,
chercher naïvement dans toutes les directions représente un travail considérable.
Ce projet s'intéresse à cette difficulté et à la manière de la surmonter.

\medskip
\noindent
À partir d'une carte réelle de Paris, nous construisons et comparons toute une
gamme de méthodes de calcul d'itinéraire, des plus élémentaires aux plus
avancées, en mesurant précisément le travail que chacune accomplit pour trouver
la même route optimale. Autour de ces méthodes, nous menons plusieurs
expériences et proposons des visualisations concrètes : la comparaison des
régions explorées sur un vrai trajet parisien, un itinéraire détaillé, une carte
interactive et une animation.

\medskip
\noindent
Les résultats montrent qu'un choix judicieux de méthode réduit le travail de
plusieurs ordres de grandeur sans jamais dégrader la qualité de la route, et
surtout que cet avantage s'accroît à mesure que la carte grandit. Ce projet
illustre ainsi, de bout en bout et sur des données réelles, une idée simple mais
puissante : face à un problème de grande taille, ce n'est pas la force brute qui
compte, mais la manière dont on prépare et oriente la recherche.
